This chapter records the special values of Dirichlet $L$-functions needed for
the analytic formula for $\hminus$.

\section{The value at \texorpdfstring{$s=0$}{s=0}}

For an odd Dirichlet character $\chi$ modulo $p$ one has the finite Hurwitz
expansion
\[
  L(s,\chi)
  = \sum_{a\in \Z/p\Z}\chi(a)\,\zeta_H^{-}\!\left(\frac{a}{p},s\right),
\]
where \(\zeta_H^{-}(x,s)\) denotes the odd Hurwitz zeta function. The
functional equation for \(\zeta_H^{-}\), evaluated at $s=0$, gives
\[
  \zeta_H^{-}(x,0) = \frac12 - x
  \qquad (0<x<1).
\]
Substituting this identity immediately yields the desired special value.

\begin{theorem}[Special value at zero]
  \label{thm:L0-eq-neg-B1}
  \lean{KummerCriterion.odd_LFunction_zero_eq_neg_BernoulliGen_one}
  \leanok
  \uses{lem:gen-bernoulli-one-intermediate, prop:gen-bernoulli-zero}
Let $p$ be an odd prime and let $\chi$ be a non-trivial odd Dirichlet
character modulo $p$. Then
\[
  L(0,\chi) = -\Bgen{1}{\chi}.
\]
\end{theorem}

\begin{proof}\leanok
Using the Hurwitz expansion and the value of \(\zeta_H^{-}(x,0)\), we obtain
\[
  L(0,\chi)
  = \sum_{a\in \Z/p\Z}\chi(a)\left(\frac12-\frac{a}{p}\right)
  = \frac12\sum_a \chi(a) - \sum_a \chi(a)\frac{a}{p}.
\]
The first sum vanishes because $\chi$ is non-trivial. The second sum is the
expression for $\Bgen{1}{\chi}$ given by
\cref{lem:gen-bernoulli-one-intermediate}. Hence
\(
  L(0,\chi) = -\Bgen{1}{\chi}.
\)
\end{proof}

\section{The value at \texorpdfstring{$s=1$}{s=1}}

For a primitive odd character $\chi$ modulo $p$, define the completed
$L$-function by
\[
  \Lambda(s,\chi)
  := \left(\frac{p}{\pi}\right)^{(s+1)/2}
     \Gamma\!\left(\frac{s+1}{2}\right)L(s,\chi).
\]
Its functional equation is
\[
  \Lambda(1-s,\chi)
  = \frac{\tau(\chi)}{i\sqrt p}\,\Lambda(s,\chi^{-1}),
\]
where
\[
  \tau(\chi)=\sum_{a\in \Z/p\Z}\chi(a)e^{2\pi i a/p}
\]
is the Gauss sum of $\chi$.

\begin{proposition}\label{thm:fe-reduction}
  \lean{KummerCriterion.odd_LFunction_one_eq_oddLValueRhs_of_LFunction_inv_zero}
  \leanok
  \uses{def:gauss-sum, def:gen-bernoulli, def:even-odd}
Let $p$ be an odd prime and let $\chi$ be a primitive odd Dirichlet
character modulo $p$. If
\(
  L(0,\chi^{-1}) = -\Bgen{1}{\chi^{-1}},
\)
then
\[
  L(1,\chi) = \frac{\pi i\,\tau(\chi)}{p}\,\Bgen{1}{\chi^{-1}}.
\]
\end{proposition}

\begin{proof}\leanok
Set $s=0$ in the functional equation. Since
\[
  \Lambda(1,\chi)=\frac{p}{\pi}L(1,\chi)
\]
and
\[
  \Lambda(0,\chi^{-1})
  = \left(\frac{p}{\pi}\right)^{1/2}\Gamma\!\left(\frac12\right)L(0,\chi^{-1})
  = \sqrt p\,L(0,\chi^{-1}),
\]
we get
\[
  \frac{p}{\pi}L(1,\chi)
  = \frac{\tau(\chi)}{i\sqrt p}\cdot \sqrt p\,L(0,\chi^{-1})
  = -i\,\tau(\chi)L(0,\chi^{-1}).
\]
Substituting the hypothesis
\(
  L(0,\chi^{-1})=-\Bgen{1}{\chi^{-1}}
\)
gives the stated formula.
\end{proof}

\begin{corollary}[Odd character formula at \texorpdfstring{$s=1$}{s=1}]
  \label{cor:L1-odd}
  \lean{KummerCriterion.odd_LFunction_one_eq_oddLValueRhs}
  \leanok
  \uses{thm:L0-eq-neg-B1, thm:fe-reduction}
Let $p$ be an odd prime and let $\chi$ be a non-trivial odd Dirichlet
character modulo $p$. Then
\[
  L(1,\chi) = \frac{\pi i\,\tau(\chi)}{p}\,\Bgen{1}{\chi^{-1}}.
\]
\end{corollary}

\begin{proof}\leanok
Because $p$ is prime, every non-trivial character modulo $p$ is primitive.
Applying \cref{thm:fe-reduction} and then
\cref{thm:L0-eq-neg-B1} to $\chi^{-1}$ gives the result.
\end{proof}

This is the form of the special-value formula needed in the proof of the
analytic class number formula for $\hminus$.
